\chapter{Resilient Execution: Collection, Failures, Sweeps and Ensembles}
\label{ch:resilient}

Part~II of the course (Lectures~13--24) turns from building graphs to running
them well. Earlier chapters have treated staging, checkpointing, native tasks,
asynchronous runs and meta-workflows (\Cref{ch:lessons-05-07,ch:meta-async,ch:native-tasks}).
This chapter covers the remaining lectures of Part~II: the scratch garbage
collector (Lecture~14), failure propagation (Lecture~16), native bindings
(Lecture~19), workflows generated from parameters (Lectures~20--21), dry runs
and portable emulation (Lecture~23), and the resumable ensemble that closes
the part (Lecture~24). Lecture~24 also exposes the one trap of name-based
checkpoints.

\section{Reference Counting and Garbage Collection}

Every \texttt{workflow://} reference increments the reference count of its
producer. The count is decremented after each consumer has executed. When it
reaches zero and the configuration sets \texttt{remove\_dir = True}, the
producer's folder is handed to the garbage collector.

\lstinputlisting[caption={Lecture 14: reference counting and collection (\texttt{examples/14-garbage/garbage\_collection.py}).},label={lst:garbage}]{../../examples/14-garbage/garbage_collection.py}

\lstinputlisting[style=bash,caption={Recorded output of Lecture 14.}]{../../examples/14-garbage/output/garbage_collection.txt}

\begin{cautionbox}[Collection renames, it does not delete]
The collector renames a folder to \texttt{<name>-removed} and records the new
path in the checkpoint. Nothing is deleted, so a mistake can be undone. The
cost is that disk space is freed only when someone removes the
\texttt{*-removed} folders after checking the results. Sinks, tasks without
consumers, are collected too.
\end{cautionbox}

\section{Failures and What They Stop}

A failed task stops exactly the part of the graph that depends on it.
Independent branches finish, and a resumed run repeats only what did not
complete.

\lstinputlisting[caption={Lecture 16: a failure and the resumed run (\texttt{examples/16-failures/fail\_and\_resume.py}).},label={lst:failures}]{../../examples/16-failures/fail_and_resume.py}

\lstinputlisting[style=bash,caption={Recorded output of Lecture 16.}]{../../examples/16-failures/output/fail_and_resume.txt}

In the first run \texttt{qc} fails, and \texttt{plot}, which depends on it,
is never started. \texttt{archive}, which depends only on \texttt{fetch},
finishes. The second run resumes from the checkpoint: \texttt{fetch} and
\texttt{archive} are reused, and only \texttt{qc} and \texttt{plot} execute.
This is the local-failure property that later chapters rely on for IoT
actuations and agentic tasks (\Cref{ch:faas-iot-agentic}).

\section{Native Bindings}

Native tasks (\Cref{ch:native-tasks}) bind workflow files and scalar values
to the parameters of an importable Python function. Lecture~19 shows three
properties of the binding: scalar values and environment variables reach the
function; an exception inside the function fails the task and its traceback is
kept in \texttt{.dagon/native\_stderr.txt}; and bindings that do not match the
function's signature are rejected when the task is constructed, before
anything runs.

\lstinputlisting[style=bash,caption={Recorded output of Lecture 19 (\texttt{examples/19-native-bindings/native\_bindings.py}).}]{../../examples/19-native-bindings/output/native_bindings.txt}

\section{Workflows Generated from Parameters}

A parameter sweep is a loop that builds tasks. Lecture~20 generates one summary task per
station and a \texttt{merge} task whose references are produced by the same loop.

\lstinputlisting[caption={Lecture 20: a sweep over stations (\texttt{examples/20-sweep/parameter\_sweep.py}).},label={lst:sweep}]{../../examples/20-sweep/parameter_sweep.py}

\lstinputlisting[style=bash,caption={Recorded output of Lecture 20.}]{../../examples/20-sweep/output/parameter_sweep.txt}

An ensemble is a sweep over perturbations of one input. Lecture~21 builds five
members whose temperatures are offset by $-1.0, -0.5, 0, 0.5, 1.0$\,\textdegree C
and a statistics task that waits for all of them.

\lstinputlisting[style=bash,caption={Recorded output of Lecture 21 (\texttt{examples/21-ensemble/ensemble.py}).}]{../../examples/21-ensemble/output/ensemble.txt}

The ensemble mean, 16.0\,\textdegree C, and the spread, 0.707\,\textdegree C
(the population standard deviation of the five member means), can both be
checked by hand.

\section{Dry Runs and Portable Emulation}

Two workflow modes let a graph be exercised without its real cost. In a
\emph{dry run} (\texttt{set\_dry}) launchers are written but not executed, so
the folders and scripts can be inspected. Under \emph{portable emulation}
(\texttt{portable\_emulation=True}) scheduler and container tasks run their
portable command locally with \texttt{bash -c}. References are rewritten to
the producers' folders without staging. LLM, FaaS and IoT tasks use their
emulated or mock providers.

\lstinputlisting[style=bash,caption={Recorded output of Lecture 23 (\texttt{examples/23-dry-portable/dry\_and\_portable.py}).}]{../../examples/23-dry-portable/output/dry_and_portable.txt}

\section{A Resumable Ensemble and the Stale-Name Trap}

Lecture~24 runs the ensemble with a checkpoint file three times, then once
more with a fresh checkpoint.

\lstinputlisting[caption={Lecture 24: a resumable ensemble (\texttt{examples/24-resumable-ensemble/resumable\_ensemble.py}).},label={lst:resumable}]{../../examples/24-resumable-ensemble/resumable_ensemble.py}

\lstinputlisting[style=bash,caption={Recorded output of Lecture 24.}]{../../examples/24-resumable-ensemble/output/resumable_ensemble.txt}

\begin{cautionbox}[Name-based reuse]
The third run doubles every offset but keeps the task names. Batch and native
checkpoints are keyed by \texttt{<workflow>.<task>} and are reused when the
recorded exit code is 0 and the folder still exists. They do not check whether
the command changed, so the third run reports the old spread, 0.707. Only a
fresh checkpoint gives the correct value, 1.414. The rule to remember is that
\emph{a changed command needs a new name or a new checkpoint}.
FaaS, IoT and agentic tasks avoid this trap by checking content: a
specification hash, output checksums, or a full identity digest
(\Cref{ch:faas-iot-agentic}).
\end{cautionbox}

\section{Chapter Summary}

\begin{itemize}[leftmargin=1.4em]
  \item Reference counts decide when a producer's folder may be collected,
    and collection renames rather than deletes.
  \item A failure stops only its descendants, and a resume repeats only
    unfinished work.
  \item Sweeps and ensembles are loops that build graphs, and their fan-in
    tasks are ordinary consumers.
  \item Dry runs and portable emulation rehearse a workflow without its cost.
  \item Name-based checkpoints are fast and simple but blind to changed
    commands.
\end{itemize}

\begin{exercise}
Modify the resumable ensemble so that each member's task name contains its
offset (for example \texttt{member\_m1.0}). Run the three-run sequence again
and explain why the trap disappears.
\end{exercise}
